\begin{table*}[!htbp]
\centering
\scriptsize
\caption{Three-term decomposition, reported separately}
\label{tab:b1}
\par\noindent\begin{minipage}{\linewidth}\raggedright{\scriptsize The decomposition is not a theorem. The terms are along observed histories, not a uniform sup over beliefs; the third term is a diagnostic refinement for observation-type policies, not an independent bound term. The phases are anchor = last informative step before the freeze, freeze = pre+1..post-1, and post = first OPEN after the freeze.}\end{minipage}\par

\par\textbf{Panel A. $W_1$ estimation term (identified posterior vs naive delta; oracle uses latent $S_t$)}

\begin{adjustbox}{max width=\linewidth}
\begin{tabular}{@{}llllllll@{}}
\toprule
$\tau$ (h) & Phase & $n$ & $W_1$ mean & $W_1$ SD & $W_1$ p90 & Oracle $W_1$ mean & Ratio \\
\midrule
$\tau=2$ & anchor & 128 & 3.96e-04 & 1.63e-19 & 3.96e-04 & 4.04e-04 & 0.980 \\
$\tau=2$ & freeze & 128 & 0.0076 & 0.0028 & 0.0105 & 0.0077 & 0.987 \\
$\tau=2$ & post & 128 & 3.96e-04 & 5.44e-07 & 3.97e-04 & 4.10e-04 & 0.966 \\
$\tau=81$ & anchor & 128 & 0.0027 & 7.56e-05 & 0.0028 & 0.0028 & 0.953 \\
$\tau=81$ & freeze & 128 & 0.0407 & 0.0160 & 0.0567 & 0.0406 & 1.003 \\
$\tau=81$ & post & 128 & 3.96e-04 & 2.70e-08 & 3.96e-04 & 4.09e-04 & 0.968 \\
\bottomrule
\end{tabular}
\end{adjustbox}
\vspace{0.5em}
\par\noindent\begin{minipage}{\linewidth}\raggedright{\scriptsize Notes. $W_1 = \mathbb{E}|\log S_t - \log P_t|$ under the identified posterior, i.e., $W_1$ between the identified posterior and the delta at $\log P_t$ (the naive treats $P_t$ as $S_t$); the oracle replaces $P_t$ by the simulator's latent $S_t$ for validation only. Ratio = identified mean / oracle mean. $n$ is the number of episodes (128 per cell); the $\tau=2$ anchor cell is numerically exact (SD $<10^{-15}$); the $\tau=81$ anchor SD is $7.6\times10^{-5}$. The freeze term grows with the freeze duration: 0.0076 at $\tau=2$ and 0.0407 at $\tau=81$.}\end{minipage}\par

\par\textbf{Panel B. Predictive-distribution TV information term}

\begin{adjustbox}{max width=\linewidth}
\begin{tabular}{@{}llllllll@{}}
\toprule
$\tau$ (h) & Phase & $n$ & TV mean & TV SD & TV p90 & TV max & Step frac TV$\ge$0.99 \\
\midrule
$\tau=2$ & anchor & 128 & 1.000 & 0.000 & 1.000 & 1.000 & 1.00 \\
$\tau=2$ & freeze & 128 & 0.646 & 0.039 & 0.686 & 0.721 & 0.50 \\
$\tau=81$ & anchor & 128 & 0.739 & 0.261 & 1.000 & 1.000 & 0.27 \\
$\tau=81$ & freeze & 128 & 0.766 & 0.066 & 0.838 & 0.911 & 0.00 \\
\bottomrule
\end{tabular}
\end{adjustbox}
\vspace{0.5em}
\par\noindent\begin{minipage}{\linewidth}\raggedright{\scriptsize Notes. TV between the identified and naive one-step predictives of $\log P_{t+1}$, evaluated at the observed history. \texttt{Step frac} is the fraction of pooled steps with TV $\ge 0.99$; episode-level maxima can be below 1 even when some steps are at 1. The post phase is not reported in the current JSON. At the $\tau=2$ anchor the two predictives are maximally different (TV = 1); during the freeze TV is 0.65 ($\tau=2$) and 0.77 ($\tau=81$).}\end{minipage}\par

\par\textbf{Panel C. Policy boundary/action-flip diagnostic}

\par\textbf{Panel C.1 Flip probability}

\begin{adjustbox}{max width=\linewidth}
\begin{tabular}{@{}llllllllll@{}}
\toprule
Policy & $\tau$ (h) & $n$ & Flip mean & Flip SD & Flip p90 & Flip max & Mean $p_{\rm id}$ & Mean $p_{\rm naive}$ & Frac flip \textgreater{} 0.5 \\
\midrule
Wide 2\% & $\tau=2$ & 128 & 0.131 & 0.081 & 0.248 & 0.329 & 0.201 & 0.073 & 0.000 \\
Wide 2\% & $\tau=81$ & 128 & 0.456 & 0.239 & 0.737 & 0.814 & 0.737 & 0.361 & 0.531 \\
Narrow 1\% & $\tau=2$ & 128 & 0.214 & 0.092 & 0.334 & 0.433 & 0.460 & 0.247 & 0.000 \\
Narrow 1\% & $\tau=81$ & 128 & 0.318 & 0.254 & 0.700 & 0.793 & 0.865 & 0.644 & 0.328 \\
\bottomrule
\end{tabular}
\end{adjustbox}
\vspace{0.5em}
\par\textbf{Panel C.2 Boundary proximity (context)}

\begin{adjustbox}{max width=\linewidth}
\begin{tabular}{@{}llllllll@{}}
\toprule
Policy & $\tau$ (h) & $n$ & Mean margin & Frac rebalance & Frac $\vert m\vert\le0.0025$ & Frac $\vert m\vert\le0.005$ & Frac $\vert m\vert\le0.01$ \\
\midrule
Wide 2\% & $\tau=2$ & 128 & -0.0126 & 0.062 & 0.031 & 0.109 & 0.281 \\
Wide 2\% & $\tau=81$ & 128 & -0.0093 & 0.133 & 0.117 & 0.219 & 0.469 \\
Narrow 1\% & $\tau=2$ & 128 & -0.0034 & 0.234 & 0.219 & 0.453 & 0.969 \\
Narrow 1\% & $\tau=81$ & 128 & -0.0021 & 0.281 & 0.250 & 0.469 & 0.953 \\
\bottomrule
\end{tabular}
\end{adjustbox}
\vspace{0.5em}
\par\noindent\begin{minipage}{\linewidth}\raggedright{\scriptsize Notes. Flip = minimal-coupling probability that the rebalance decision at $t+1$ differs, $|p_{\rm id}-p_{\rm naive}|$, where $p_{\rm model} = P(|P_{t+1}/{\rm anchor}-1| > \theta)$ under that model's predictive; positions are from the observed trajectory (first-order). $p_{\rm id}$ and $p_{\rm naive}$ are the model-implied rebalance probabilities. Margin = $|P_t/{\rm anchor}-1| - \theta$ at decision steps; $m>0$ means the observed path rebalances. The boundary-proximity panel is context for the flip diagnostic, not a fourth term. The flip diagnostic is not a bound: it describes how close the policy is to its decision boundary under the two models.}\end{minipage}\par

\end{table*}
